% Unit 9 map -- one page.
\documentclass{apchem-worksheet}

\apcunit{9}
\apcunittitle{Thermodynamics and Electrochemistry}
\apcdoctitle{Unit Map}
\apcsubtitle{CED Unit 9 \textbullet\ \textbf{7--9\%} of the AP Exam \textbullet\ 9 blocks + float + exam}

\begin{document}
\apctitleblock

\begin{apcnote}[How this unit is graded --- and what makes it different]
Unit exam \textbf{Fri Mar 19}: 25 multiple choice (\textbf{no calculator})
$+$ 1 short and 1 long free-response, 69 minutes, 39 points.

This is the last unit of the course, and it is the one that \emph{closes}
it. Unit 6 asked whether a reaction releases heat. Unit 7 asked how far it
goes. Unit 5 asked how fast. Unit 9 supplies the single quantity that ties
all three together --- and then builds a machine out of it.

The whole unit runs on one chain of equalities:
\[ \Delta G^\circ = \Delta H^\circ - T\Delta S^\circ
   \qquad \Delta G^\circ = -RT\ln K
   \qquad \Delta G^\circ = -nFE^\circ_{\text{cell}} \]
Three ways of writing the same fact. A negative $\Delta G^\circ$, a $K$
greater than 1, and a positive cell voltage are one statement in three
languages, and most of this unit is translation between them.
\end{apcnote}

\section*{\sffamily Block calendar}

\renewcommand{\arraystretch}{1.25}
\noindent\begin{tabular}{@{}p{2.1cm}p{4.9cm}p{1.6cm}p{3.9cm}p{1.9cm}@{}}
\toprule
\textbf{Date} & \textbf{Topics} & \textbf{CED} & \textbf{Read (PDF pp.)} & \textbf{Practice}\\
\midrule
Mon Feb 22 & Entropy; predicting the sign of $\Delta S$ & 9.1 & \S17.1 \textbullet\ 835--841 & WS 1\\
Wed Feb 24 & Absolute entropy; $\Delta S^\circ$ from tables & 9.2 & \S17.2, 17.5--17.6 \textbullet\ 841--853 & WS 1\\
Fri Feb 26 & Gibbs free energy; the four cases & 9.3 & \S17.3--17.4 \textbullet\ 842--848 & WS 2\\
Mon Mar 1 & $\Delta G_f^\circ$ route; thermodynamic vs kinetic control & 9.3, 9.4 & \S17.4, 17.7 \textbullet\ 845--859 & WS 2\\
Wed Mar 3 & Free energy and equilibrium; $\Delta G^\circ = -RT\ln K$ & 9.5 & \S17.9 \textbullet\ 862--867 & WS 3\\
Fri Mar 5 & Free energy of dissolution; coupled reactions & 9.6, 9.7 & \S17.5, 17.10 \textbullet\ 848--849, 867--869 & WS 3\\
Mon Mar 8 & Galvanic and electrolytic cells & 9.8 & \S18.1, 18.7 \textbullet\ 882--884, 906--911 & WS 4\\
Wed Mar 10 & Cell potential and free energy; $E^\circ$, $\Delta G^\circ$, $K$ & 9.9 & \S18.2--18.3 \textbullet\ 884--895 & WS 4\\
Fri Mar 12 & Nonstandard conditions; electrolysis and Faraday & 9.10, 9.11 & \S18.4, 18.7 \textbullet\ 895--900, 906--911 & WS 5 (FRQ)\\
Mon Mar 15 & Reteach / float & --- & review sheet & ---\\
Wed Mar 17 & Unit review & all & review sheet & ---\\
\textbf{Fri Mar 19} & \textbf{UNIT 9 EXAM} & all & review sheet & ---\\
\bottomrule
\end{tabular}

{\sffamily\footnotesize Dates from the co-op calendar. This is the final
content unit; cumulative review begins Mon Mar 29 after the Easter break.}

\vspace{0.4em}
\section*{\sffamily By the end of this unit you can\ldots}

\begin{itemize}[leftmargin=*,itemsep=0.1em]
  \item Predict the sign of $\Delta S^\circ$ from the physical change.
        \ced{9.1}
  \item Calculate $\Delta S^\circ$ from absolute entropies. \ced{9.2}
  \item Calculate $\Delta G^\circ$ two ways and judge thermodynamic
        favorability. \ced{9.3}
  \item Explain why a favored reaction may not proceed at a measurable
        rate. \ced{9.4}
  \item Relate $\Delta G^\circ$ to $K$, and $\Delta G$ to $Q$. \ced{9.5}
  \item Identify the entropy and enthalpy contributions to dissolution.
        \ced{9.6}
  \item Explain how coupling or an external energy source drives an
        unfavorable process. \ced{9.7}
  \item Distinguish galvanic from electrolytic cells and label the
        electrodes. \ced{9.8}
  \item Relate $E^\circ_{\text{cell}}$ to $\Delta G^\circ$ and to $K$.
        \ced{9.9}
  \item Predict qualitatively how concentration changes a cell potential.
        \ced{9.10}
  \item Apply Faraday's law to electrolysis. \ced{9.11}
\end{itemize}

\begin{apctrap}[Where students lose points in this unit]
Mixing J and kJ in $\Delta H^\circ - T\Delta S^\circ$ (the single most
common arithmetic error in the unit); setting $S^\circ$ of an element to
zero; multiplying $E^\circ$ when a half-reaction is scaled; taking $n$ from
a half-reaction instead of the balanced overall equation; forgetting to
convert time to seconds in Faraday's law; and calling a favored reaction
``fast''.

Two conventions the CED is strict about. Say
\textbf{thermodynamically favored}, never ``spontaneous'' --- Zumdahl's
chapter title still uses the old word. And topic 9.8 \textbf{excludes}
labelling electrodes as positive or negative; identify them as
\emph{anode} (oxidation) and \emph{cathode} (reduction) instead.
\end{apctrap}

\end{document}
